\section{Examples}

\subsection{RAG Workflow}

A baseline enterprise support workflow may retrieve 20 document chunks, send all chunks to a frontier model, generate an answer, and return it. A \migaki-optimized plan may deduplicate overlapping chunks, drop low-relevance chunks, preserve stable policy instructions as a cacheable prefix, route chunk ranking to a cheaper model, use a stronger model for synthesis, validate the answer against retrieved sources, retry only the synthesis step if validation fails, and emit an evidence bundle.

The agent framework still decided the strategy:

\begin{center}
Retrieve $\rightarrow$ Answer $\rightarrow$ Validate.
\end{center}

\migaki{} optimized the execution:

\begin{center}
Deduplicate $\rightarrow$ Cache $\rightarrow$ Route $\rightarrow$ Parallelize $\rightarrow$ Validate $\rightarrow$ Report.
\end{center}

\subsection{Code Review Workflow}

A code-review agent may load repository context, changed files, style guides, and prior conversation, then ask a frontier model for review. A \migaki{} plan can mark style guides as fixed and cacheable, mark changed files as non-droppable, remove unrelated history, run deterministic static checks, route lint explanations to a cheaper model, use a stronger model for architectural review, validate final comments against changed lines, and export evidence.

This makes execution policy explicit and measurable: baseline tokens, optimized tokens, baseline cost, optimized cost, latency, validator pass rate, false-positive rate, and comment acceptance rate.
